\chapter{项目概述}

\section{研究背景与意义}
随着微波通信技术的快速发展，电磁污染日益严重，不仅干扰电子设备，也对人体健康构成威胁。微波吸收材料（MAMs）被视为一种有效
的解决方案，它们能将微波能量转化为其他形式，例如热能。目前已探索了多种 MAMs，包括一维材料（如 SiC 纳米线、碳纳米管等）、
二维材料（如石墨烯、MXene 等）、以及块体材料（如碳基复合材料、陶瓷等）。然而，许多 MAMs 仍存在有效吸收带宽窄、密度过大的问
题，限制了其满足实际应用中“薄厚度、低密度、强吸收、宽频带”的微波吸收需求。

在这些候选材料中，聚合物衍生陶瓷（PDCs）因其内在优势，近年来受到了广泛关注，包括分子水平的前驱体均一性、相较于传统陶瓷
路线相对较低的加工温度，以及设计具有定制功能的新型成分的可行性。然而，传统的聚合物衍生陶瓷通常表现出有限的吸收带宽和较差的阻抗匹配，这导致入射
微波被显著反射，最终限制了它们的实际应用。

最近的研究表明，对聚合物衍生陶瓷进行成分调控可以有效调节电磁参数并增强微波吸收。例如，Li 等人\cite{li2025electromagnetic} 通过聚合
物衍生陶瓷法制备了含有 Fe$_3$Si 球和 SiO$_2$ 纳米线的 SiFeBOC 陶瓷，在仅 1.6~mm 厚度下实现了 $RL_{\min} = -47.68$~dB 的最小反射损耗
和 4.39~GHz 的有效吸收带宽，这归功于磁电耦合、界面极化以及改善的阻抗匹配的协同效应。Chen 等人\cite{CHEN2025181401} 在 Sm 掺杂 SiBCN 陶瓷中
提出了一种可控的相变策略，通过波透 Si$_3$N$_4$、有损 SiC 和石墨的联合效应，实现了 $RL_{\min} = -54.24$~dB 和 7.48~GHz 的宽带宽。
\section{国内外研究现状}
近年来，国际上围绕硼化物、碳化物与氧化物复合体系开展了大量研究，重点关注高温结构稳定性与辐射特性耦合调控\cite{wang2023hightemp,liu2024emissivity}。
Zhao等\cite{zhao2022a}通过在碳化硅材料表面镀银,形成导电网络通路,提高其电磁损耗;增强界面极化,提高其衰减常数和阻抗匹配, 从而增强其吸波性能。
Liu等\cite{liu2025}通过超快高温烧结的方式制备了三种新型HEP材料，通过调整组分并触发多种电磁损耗机制，该体系的电磁吸收性能得到了显著提升。
Zhou等\cite{zhou2019}通过化学气相渗透法构建了SiC$_f$/Si$_3$N$_4$复合材料,实现了在 宽温度范围(25–600°C)内X和Ku波段对温度不敏感的电磁波吸收。
Lan和Wang等\cite{lan2020}通过在碳化硅晶须组成的轻质多孔骨架中掺入碳纳米线，热还原以实现碳化硅纳米线网络以产生强界面极化和多重反射，提高电导率。
Cai等\cite{cai2024}通过在1500摄氏度、氮气气氛下烧结碳化硅晶须的方式实现了碳化硅/氮化硅复合气凝胶，将Si$_3$N$_4$纳米带引入碳化硅纳米线网络,增强了阻抗匹配、界面极化和偶极极化。
我们总结了学界的部分代表性工作，见表~\ref{tab:EAB}，以供参考。
\input{maintext/EAB}
\section{研究目标}
本项目立项时的核心目标如下：旨在开发一种具有非晶微观结构的新型 SiZrNOC 纳米纤维膜（NFM），以解决传统陶瓷纤维在高温环境下隔热性能不足的问题。
团队通过：
\begin{itemize}
    \item 引入 Zr 元素；
    \item 优化纤维的微观结构，提升其热稳定性和机械性能。
\end{itemize}
从而满足航空航天等领域对高性能隔热材料的需求。